\begin{figure}[H]
  \centering
  \includegraphics[width=\linewidth]{figures/figure4.pdf}
  \caption{\new{\textbf{Inter-annotator agreement and hallucination analysis.} (a)~Mention-level inter-annotator agreement per cross-annotated note (F1-based IAA, Cohen's $\kappa$, and PABAK); pooled across the cross-annotated subset, F1$=0.837$, $\kappa=0.597$, PABAK$=0.609$. (b)~Pre-specified structural-disambiguation rule triggers ($n=552$; 15.6\% of candidate-mention rows), the principal source of apparent annotator divergence. (c)~Seven-class hallucination taxonomy applied to candidate false positives, with population-level rates extrapolated to the full false-positive set ($N=7{,}171$) by post-stratified Horvitz--Thompson estimation. (d)~Annotation gap versus true hallucination: 54.3\% of apparent false positives are not true errors (annotation gap, synonym equivalence, or guideline ambiguity), leaving a 45.7\% true-hallucination share of false positives ($\approx 24.4\%$ of all predictions).}}
  \label{fig:iaa-halluc}
\end{figure}
